% Tableau des labs du cours HE2 (SEC4062) dans NDAP SUKULU (propre au cours).
Le cours compte dix laboratoires, disponibles sous le code \codeUE. Ils
travaillent tous sur des pièces (relevés, journaux, exports, résultats de
balayage) et jamais sur un système réel~: ils entraînent le raisonnement du
testeur et du défenseur, pas le geste d'attaque.

\begin{tableaucours}{L{2.4cm}L{6.4cm}L{1.6cm}L{1.8cm}}{%
    \tableauHeader{Lab} &
    \tableauHeader{Sujet} &
    \tableauHeader{Chapitre} &
    \tableauHeader{Opération}
}
HE2-LAB-01 & Périmètre~: trier les cibles autorisées & \ref{chap:engagement} &
LOGONE \\
HE2-LAB-02 & OSINT~: sous-domaines et exposition & \ref{chap:osint} & LOGONE \\
HE2-LAB-03 & Lire un résultat de balayage & \ref{chap:balayage} & LOGONE \\
HE2-LAB-04 & Prioriser des vulnérabilités & \ref{chap:vulnerabilites} &
LOGONE \\
HE2-LAB-05 & Active Directory~: chemins vers les administrateurs &
\ref{chap:ad} & LOGONE \\
HE2-LAB-06 & Élévation de privilèges Linux~: relevé & \ref{chap:privileges}
& LOGONE \\
HE2-LAB-07 & Audit de mots de passe & \ref{chap:mdp} & LOGONE \\
HE2-LAB-08 & Historique HTTP~: contrôle d'accès et injection &
\ref{chap:injections} & MEFOU \\
HE2-LAB-09 & Équipe violette~: couverture de détection & \ref{chap:rouge} &
MBAM \\
HE2-LAB-10 & Rapport~: calculer les scores CVSS & \ref{chap:rapport} & LOGONE \\
\end{tableaucours}
\vspace{-0.8em}
\begin{center}{\small\itshape Les laboratoires du cours dans NDAP SUKULU}\end{center}
